\section{Connectivity and defect}\label{sec:defect}

\begin{lemma}\label{lem:eight-complement}
Every eight-vertex graph $J$ of minimum degree at least five contains a
$K_6^{-}$ minor.
\end{lemma}

\begin{proof}
The complement of $J$ has maximum degree at most two.  Extend it to an
edge-maximal graph $F$ with that property.  Every component of $F$ is a
cycle, except possibly one component which is $K_1$ or $K_2$: any longer
path could be closed, and deficient vertices in two components could be
joined.  The seven possibilities for $F$ and explicit models in its
complement are displayed below.  Cycle labels are consecutive; a string
such as $03$ denotes the set $\{0,3\}$.
\[
\begin{array}{c|c|c}
F&\text{six branch sets}&\text{only missing pair}\\\hline
C_8(0,\ldots,7)&03,15,2,4,6,7&6:7\\
C_5(0,\ldots,4)\dot\cup C_3(5,6,7)&05,1,26,3,4,7&3:4\\
C_4(0,\ldots,3)\dot\cup C_4(4,\ldots,7)&04,1,25,3,6,7&6:7\\
C_7(0,\ldots,6)\dot\cup K_1(7)&024,1,3,5,6,7&5:6\\
C_4(0,\ldots,3)\dot\cup C_3(4,5,6)\dot\cup K_1(7)
 &024,1,3,5,6,7&5:6\\
C_6(0,\ldots,5)\dot\cup K_2(6,7)&024,1,3,5,6,7&6:7\\
C_3(0,1,2)\dot\cup C_3(3,4,5)\dot\cup K_2(6,7)
 &03,14,2,5,6,7&6:7
\end{array}
\]
Each nonsingleton is connected in $\overline F$, and all interbag
adjacencies except the indicated one are present.  Since $\overline F$
is a spanning subgraph of $J$, each displayed model also lies in $J$.
\end{proof}

Theorem~\ref{thm:defect-ladder} needs an edge of codegree at most three
without a colouring hypothesis.  Contracting exterior components reduces
the degree-seven and degree-eight cases to the following finite lemma.

\begin{lemma}[Computer-assisted finite lemma]\label{lem:finite-quotients}
For a graph $J$ and $A\subseteq V(J)$, let $Q(J,A)$ be obtained by adding
nonadjacent vertices $v,c$, making $v$ complete to $J$ and giving $c$
precisely the neighbourhood $A$.
\begin{enumerate}
\item If $|V(J)|=7$, $\delta(J)\geq4$ and $|A|\geq6$, then
$K_7^{-}\minor Q(J,A)$.
\item Fix a labelled graph $J$ with $|V(J)|=8$, $\delta(J)\geq4$ and
$K_6^{-}\not\minor J$.  At most one set $A\subseteq V(J)$ with $|A|\geq6$
makes $Q(J,A)$ $K_7^{-}$-minor-free.  If such a set exists,
$V(J)\setminus A$ consists of two adjacent vertices of degree four in $J$.
\end{enumerate}
\end{lemma}

\begin{proof}[Finite verification and its scope]
For part (1), $\overline J$ has maximum degree two, hence is a union of
paths and cycles.  Enumerating the multisets of such components with total
order seven gives 29 types.  For each, check the full attachment and its
seven possible single omissions.  Each of the $29\cdot8=232$ nine-vertex
quotients has a seven-bag model.  The checker enumerates all 750 partitions
of vertex subsets into seven nonempty bags, and tests connectivity and all
interbag adjacencies.  A seven-bag model on nine vertices uses seven, eight
or nine vertices, so this search is exhaustive.

For part (2), adjoin a vertex with each possible neighbourhood to each of
the 1,044 unlabelled seven-vertex graphs in the NetworkX graph atlas.
Every eight-vertex graph occurs.  The minimum-degree filter and exact
isomorphism testing leave 424 classes.
Of these, 55 have no $K_6^{-}$ minor.  Checking all missed sets of order
zero, one or two gives $55(1+8+28)=2,035$ exterior profiles.  Exactly 2,031
have certified $K_7^{-}$ models.  Each of the four remaining profiles
has the uniqueness and degree properties in part (2).

Here the minor engine starts with singleton bags and recursively either
deletes a bag or merges two touching bags.  These operations preserve
disjoint connected bags and generate every minor model by contracting
spanning trees and deleting unused vertices.  At six or seven bags it
accepts precisely when at most one pair is nonadjacent.  Every positive
certificate is separately checked for disjointness, connectivity and
interbag adjacencies.  The implementations test known positive and negative
instances; a separate internal audit also checked the four negative
quotients by all 11,880 seven-bag partitions of their vertex subsets.

The accompanying reproducibility note, \texttt{README.md}, records the
four profiles, executable inputs and SHA-256 certificate digests.
Part (1) uses only the Python standard library; part (2) also uses
NetworkX's atlas and exact isomorphism implementation.  The reductions
below impose no bound on the order of an exterior component or of $G$.
\end{proof}

\begin{theorem}\label{thm:six-connected-eight}
In every six-connected $K_7^{-}$-minor-free graph, each degree-eight vertex
is incident with an edge of codegree at most three.
\end{theorem}

\begin{proof}
Let $v$ have degree eight and put $J=G[N(v)]$.  Suppose every incident
codegree is at least four.  Then $\delta(J)\geq4$, and $J$ has no
$K_6^{-}$ minor because adjoining $\{v\}$ would give the forbidden
minor.  If $G-N[v]$ is empty, minimum degree six in $G$ gives
$\delta(J)\geq5$, contradicting Lemma~\ref{lem:eight-complement}.

Every component $C$ of $G-N[v]$ has at least six neighbours in $N(v)$,
since its neighbourhood separates it from $v$.  Contracting $C$ to $c$
and deleting the other exterior components gives exactly
$Q(J,N_G(C))$.  Lemma~\ref{lem:finite-quotients}(2) therefore applies to
every exterior component.  For the fixed labelled graph $J$, all
components must miss the same unique pair $a,b$.  Neither $a$ nor $b$
has an exterior neighbour, so both have degree $1+4=5$ in $G$,
contradicting six-connectivity.
\end{proof}

The degree-six case adapts a computation-free argument of Norin and
Totschnig.  We verify the hypotheses under $K_7^{-}$ exclusion.

\begin{lemma}[Degree-six disc bound]\label{lem:degree-six}
Let $H$ be a five-connected $K_7^{-}$-minor-free graph on at least nine
vertices.  If every edge has codegree at least four and some vertex has
degree six, then $|E(H)|\leq4|V(H)|-9$.
\end{lemma}

\begin{proof}
Let $v$ have degree six.  Since $d_{H[N(v)]}(x)=c(vx)\geq4$, label
$N(v)=\{u_1,w_1,u_2,w_2,u_3,w_3\}$ so that all its possible nonedges
are $u_iw_i$.  Two disjoint paths joining two distinct displayed pairs, with open
interiors outside $N[v]$,
would fill two of the three missing edges after contraction.  Together
with $v$ they give a $K_7^{-}$ model.  We call this the two-pair
obstruction; a present pair-edge may serve as one of the paths.

All three displayed pair-edges are absent.  Otherwise an exterior
component, with at least five neighbours in $N(v)$ by five-connectivity,
joins a second pair and violates the two-pair obstruction.  Moreover,
some pair, say $u_1,w_1$, has no exterior common neighbour.  Otherwise
choose a common neighbour for each pair.  Two distinct choices contradict
the obstruction.  If all choices coincide at $x$, order at least nine
leaves a component outside $N[v]\cup\{x\}$.  Its boundary has at least
five vertices, at least four of which belong to $N(v)$.  It therefore
joins one pair disjointly from a second pair's path through $x$, again a
contradiction.

We now use the two-paths theorem in the form of
\cite[Theorem~13]{NorinTotschnig2025}.  The next two applications are the
degree-six argument in \cite[Claims~3.12--3.15]{NorinTotschnig2025}; we
record the separator and density hypotheses needed for them.  There is
no nontrivial separation $(A,B)$ of order five with $N[v]\subseteq A$.
Indeed, five-connectivity and Menger's theorem provide five disjoint paths
in $H[A]$ from five distinct members of $N(v)$ to the five boundary
vertices, with interiors avoiding $N[v]$.  These five members contain
two complete matching pairs.  Temporarily label their boundary ends
$s_1,t_1$ and $s_2,t_2$, and label the fifth end $x$.  Apply the two-paths
theorem in $H[B]-x$ to the order $s_1,s_2,t_1,t_2$.  A crossing violates
the two-pair obstruction.  A separation outcome lifts after restoring
$x$ to a cut of order at most four in $H$.  In the remaining outcome the
graph embeds in a disc with the four terminals on its boundary.  An edge
with an end in $B\setminus A$ has at most two common neighbours in that
disc: three would enclose a vertex behind a triangle, which together
with $x$ gives a cut of order at most four in $H$.  Restoring $x$ gives
codegree at most three, a contradiction.

Apply the same theorem to $H'=H-\{v,u_1,w_1\}$, with terminal order
$u_2,u_3,w_2,w_3$.  A crossing again violates the two-pair obstruction.
A separating outcome lifts, after restoring $u_1,w_1$, to a separation
of order at most five with $N[v]$ on one side; five-connectivity and the
preceding paragraph exclude it.  Thus $H'$ embeds in a disc with the
four terminals on its boundary.  Put $n'=|V(H')|$.  Then
$|E(H')|\leq3n'-7$.  Each of its vertices outside $N[v]$ sees at most
one of $\{v,u_1,w_1\}$; the four remaining neighbourhood vertices see
all three, and the deleted triple induces two edges.  Hence
\[
 |E(H)|\leq(3n'-7)+(n'-4)+12+2=4n'+3=4|V(H)|-9.
\]
The forbidden-minor hypothesis was used only for the displayed
$K_7^{-}$ two-pair model.  Neither edge-maximality nor the stronger
$K_7^{\vee}$ exclusion is used in this argument.
\end{proof}

\begin{proposition}\label{prop:density-edge}
Let $G$ be six-connected and $K_7^{-}$-minor-free.  If
$|E(G)|\geq4|V(G)|$, then some edge of $G$ has codegree at most three.
\end{proposition}

\begin{proof}
The density implies order at least nine, so
Lemma~\ref{lem:jakobsen-defect} gives average degree less than nine.  Thus a
vertex of minimum degree has degree six, seven or eight.  Suppose every
edge has codegree at least four.  Lemma~\ref{lem:degree-six} excludes
degree six, and Theorem~\ref{thm:six-connected-eight} excludes degree
eight.  At a degree-seven vertex $v$, the local graph $J=G[N(v)]$ has
minimum degree at least four.  There is an exterior component, and
six-connectivity gives it at least six boundary neighbours.  Contracting
it produces a quotient forbidden by Lemma~\ref{lem:finite-quotients}(1).
This excludes all three degrees.
\end{proof}

\begin{proof}[Proof of Theorem~\ref{thm:defect-ladder}]
Let $G$ satisfy the hypotheses.  Proposition~\ref{prop:density-edge}
supplies an edge $e$ of codegree at most three.  The quotient $G/e$ is
$(r-1)$-connected and has no $K_7^{-}$ minor.  By \eqref{eq:contraction},
$|E(G/e)|\geq4|V(G/e)|$ and $D(G)\geq D(G/e)+1$.
The retained density ensures $|V(G/e)|\geq9$.

For $r=6$, Lemma~\ref{lem:jakobsen-defect} gives $D(G/e)\geq25$, so
$D(G)\geq26$.  For $r\geq7$, induction on $r$ applies to $G/e$ and gives
\[
              D(G)\geq D(G/e)+1\geq20+(r-1)+1=20+r.
\]
\end{proof}

The base case uses Jakobsen's bound because contraction need not preserve
six-connectivity.
